\smartsection{Image Processing}
\label{sec:image-processing}

\smartsubsection{Overview \& Architecture}

The image-processing module is responsible for turning a raw video feed
into three things the rest of the system consumes: an annotated frame for
the live dashboard, a person count for telemetry, and a debounced trigger
for e-mail/MQTT notifications (Section~\ref{sec:mqtt-email}). It runs
entirely in Python — the one part of the assignment explicitly exempted
from the "core logic must be C" requirement — and was designed around a
single constraint: the Orange Pi Zero Plus 2 has no camera input of its
own.

\smartsubsubsection{Detector}

The primary detector is OpenCV's built-in \keyterm{HOG+SVM person
detector} 

(\texttt{cv2.HOGDescriptor\_getDefaultPeopleDetector()}). This
was chosen deliberately over a deep-learning alternative: it ships with
OpenCV (no model file to download or manage), runs entirely on the CPU,
and is light enough to be workable on the board's Cortex-A53 cores without
any hardware acceleration. A secondary, heavier \texttt{mobilenet\_ssd}
backend is also implemented and selectable via configuration, for cases
where accuracy matters more than frame rate.

Beyond the baseline detector, a second module, \texttt{fused\_tracker.py},
was developed to combine the HOG detector with three Haar-cascade cues
(frontal face, profile face, upper body) as auxiliary evidence, and — most
relevantly for Section~\ref{sec:resolution-experiment} — to
\keyterm{downsample the frame before detection} via a configurable
\texttt{detection\_scale} factor, rescaling any resulting bounding boxes
back to full resolution afterward. This is the mechanism behind the
"optimized" resolution results discussed later in this section.

\smartsubsubsection{No-Camera Workaround: HTTP Polling}

Since the board has no camera, the webcam is captured on the
\textbf{laptop} and made available to the board over the network. Two more
conventional approaches were tried first and both failed in practice:

\begin{itemize}
    \item A raw \textbf{UDP} video stream between the laptop and the
    board.
    \item An \textbf{HTTP MJPEG} stream, decoded by OpenCV's FFmpeg
    backend on the receiving end.
\end{itemize}

Both consistently froze on the first frame when fed from a live network
source into \texttt{cv2.VideoCapture} — a known class of issue with
OpenCV's FFmpeg demuxer and non-seekable, live streams. The working
solution replaces both with plain \keyterm{HTTP polling}: the laptop runs
a small HTTP server (\texttt{laptop\_camera\_server.py}) that always
serves whatever the most recently captured JPEG frame is at
\texttt{GET /latest.jpg}; the board's \texttt{HttpPollCapture} class
(a drop-in replacement for \texttt{cv2.VideoCapture} within

\texttt{person\_detector.py}) simply issues a fresh HTTP GET for every
frame it needs. There is no persistent stream, no protocol state, and
therefore nothing to get stuck — each request is independent.

\begin{table}[H]
\centering
\caption{The four-process capture pipeline, spanning both machines.}
\label{tab:capture-pipeline}
\resizebox{\textwidth}{!}{%
\begin{tabular}{lll}
\toprule
\textbf{Process} & \textbf{Runs on} & \textbf{Role} \\
\midrule
\texttt{laptop\_camera\_server.py} & Laptop & Captures webcam, serves latest JPEG over HTTP \\
\texttt{person\_detector.py} & Board & Polls the laptop, runs detection, writes results \\
\texttt{frame\_viewer\_server.py} (via \texttt{pi\_stream\_result.sh}) & Board & Serves the annotated frame back out \\
\texttt{laptop\_view\_result.sh} & Laptop & Opens a browser view of the annotated feed \\
\bottomrule
\end{tabular}
}
\end{table}

\smartsubsubsection{Overlays and Shared Output}

Every processed frame is annotated with a bounding box per detection, a
live person counter, the student ID, the current system date/time, and a
smoothed FPS measurement (Figure~\ref{fig:720p-result} shows a
representative annotated frame together with the live telemetry it feeds).
The frame and the detection result are then published for the rest of the
system to consume — most importantly, the C-based web server and REST API
covered in later sections — via two files written to a RAM-backed
\texttt{tmpfs} (\texttt{/dev/shm/surveillance/}), chosen so that neither
the video frames nor the frequent JSON writes wear down the SD card:

\begin{lstlisting}[style=pythonstyle_bright_2026, caption={Person-count and telemetry payload written on every processed frame.}, label={code:persons-json}]
payload = {
    "count": count,
    "timestamp": datetime.now().isoformat(timespec="seconds"),
    "fps": round(fps, 2),
}
data = json.dumps(payload).encode("utf-8")
self._atomic_write_bytes(self.persons_path, data)
\end{lstlisting}

Both the frame and the JSON file are written atomically (temp file, then
an OS-level rename) so that any reader — including the C code in
Section~\ref{sec:web-server} — never observes a partially written file and
therefore needs no locking.

\smartsubsection{Experiment 3-1: Detection Accuracy Under Varying Lighting}

Detection accuracy was measured under four lighting conditions — daylight,
artificial (indoor) light, low light, and strong backlight — using
\texttt{detection\_logger.py} to poll the live person count while a human
observer recorded the ground-truth number of people actually in frame. A
minimum of eight samples was collected per condition, covering a mix of
distances, poses, and backgrounds (see the raw logs in
\texttt{results/light-results/} for the full per-sample breakdown).

\begin{table}[H]
\centering
\caption{Detection accuracy by lighting condition (correct = detected count exactly matches ground truth).}
\label{tab:lighting-accuracy}
\begin{tabular}{lccc}
\toprule
\textbf{Condition} & \textbf{Correct} & \textbf{Total Samples} & \textbf{Accuracy} \\
\midrule
Daylight & 10 & 11 & \stable{90.9\%} \\
Artificial light & 8 & 10 & \stable{80.0\%} \\
Low light & 6 & 11 & \highlight{54.5\%} \\
Backlight & 3 & 8 & \highlight{37.5\%} \\
\bottomrule
\end{tabular}
\end{table}

The trend is monotonic and unsurprising: accuracy degrades as the amount
and quality of usable light decreases, with \alert{backlight conditions
being the most damaging} — a strongly backlit subject is reduced to a
near-silhouette, destroying most of the gradient information the HOG
descriptor relies on. Two representative frames illustrate this:

\begin{figure}[H]
\centering
% \includegraphics[width=0.4\textwidth]{../../2.Image Processing/results/fig-gitignore/320p_1_1_artificial-light.png}
\caption{An artificial light frame in which the subject is detected nicely (good aligned) bounding box.}\label{fig:artificial-light}
\end{figure}

\begin{figure}[H]
\centering
% \includegraphics[width=0.4\textwidth]{../../2.Image Processing/results/fig-gitignore/320p_1_1_back-light.png}
\caption{A backlit frame in which the subject is reduced to a silhouette; the detector nonetheless produced a (loosely aligned) bounding box, illustrating both the difficulty of this condition and the detector's partial robustness to it.}\label{fig:back-light-hit}
\end{figure}

\begin{figure}[H]
\centering
% \includegraphics[width=0.4\textwidth]{../../2.Image Processing/results/fig-gitignore/320p_0_1_low-light.png}
\caption{A low-light frame where the detector produced no bounding box despite a clear, close frontal view — one of the accuracy failures in Table~\ref{tab:lighting-accuracy}.}\label{fig:low-light-miss}
\end{figure}

\begin{figure}[H]
\centering
% \includegraphics[width=0.4\textwidth]{../../2.Image Processing/results/fig-gitignore/320p_1_1_extreme-low-light.png}
\caption{An extreme-low-light frame where the detector could detect the subject clearly and box it nicely.}\label{fig:extreme-low-light}
\end{figure}

\begin{figure}[H]
\centering
% \includegraphics[width=0.4\textwidth]{../../2.Image Processing/results/fig-gitignore/320p_1_1_farther.png}
\caption{A farther subject detection went smoothly and the box was closely aligned with the subject.}\label{fig:farther}
\end{figure}

\begin{figure}[H]
\centering
% \includegraphics[width=0.4\textwidth]{../../2.Image Processing/results/fig-gitignore/320p_1_1_hand-test.png}
\caption{On this test we changed the subject to a body part (in this case a hand) which was successfully detected using our cascade model and the box alignment was not bad.}\label{fig:hand-test}
\end{figure}

\begin{figure}[H]
\centering
% \includegraphics[width=0.4\textwidth]{../../2.Image Processing/results/fig-gitignore/320p_2_1_back-light.png}
\caption{A strongly backlit frame where the detector failed to produce bounding box around the subject due to the high backlight — one of the accuracy failures in Table~\ref{tab:lighting-accuracy}.}\label{fig:back-light-miss}
\end{figure}

\smartsubsection{Experiment 3-2: Spoof Test}

\begin{questionsection}{Can a printed photo or phone/tablet screen fool the detector?}
The assignment requires testing the detector against a static image of a
person (printed or displayed on a screen) rather than a live subject, and
asking whether it fools the system. \textit{Guidance:} consider what
information the detector actually uses to decide "person" versus
"not person" and whether that information distinguishes a live subject
from a flat image of one.
\end{questionsection}

The test was performed by holding a tablet displaying a photograph of the
test subject in front of the camera at several angles.

\begin{figure}[H]
\centering
% \includegraphics[width=0.4\textwidth]{../../2.Image Processing/results/fig-gitignore/320_0_0_spoof-test.png}
\caption{Tablet held facing the camera, straight-on: the detector correctly reports zero persons.}\label{fig:spoof-miss}
\end{figure}

\begin{figure}[H]
\centering
% \includegraphics[width=0.4\textwidth]{../../2.Image Processing/results/fig-gitignore/320p_0_1_spoof-test.png}
\caption{Same tablet, slightly different angle/framing: the detector now reports one person and draws a bounding box over the displayed photo.}\label{fig:spoof-hit}
\end{figure}

\begin{answersection}{Findings and proposed mitigation}
The result was \punder{inconsistent rather than a clean pass or fail}: in
some framings the detector correctly ignored the tablet
(Figure~\ref{fig:spoof-miss}), while in others — evidently once the
photo's pose and scale happened to resemble the gradient pattern the HOG
descriptor was trained on — it registered a false positive
(Figure~\ref{fig:spoof-hit}). This is expected: a HOG/SVM detector (and
equally a CNN-based one without additional safeguards) has no notion of
\keyterm{liveness}. It only evaluates the 2-D gradient structure of the
image region and cannot distinguish a flat photograph from a real
three-dimensional subject.

Three mitigations were considered for future work, none of which were
implemented in the current system:
\begin{itemize}
    \item \textbf{Depth sensing} — a stereo camera or IR depth sensor
    would trivially reject any perfectly flat surface.
    \item \textbf{Motion-based liveness} — requiring small, natural
    movement across several consecutive frames before confirming a
    detection would defeat a static photo, though not a looping video.
    \item \textbf{Texture/reflection analysis} — screens and glossy
    prints reflect specular light differently than skin and fabric; a
    lightweight secondary classifier on the detected region could flag
    this.
\end{itemize}
Given the board's limited compute budget, motion-based liveness is the
most realistic near-term addition, since it requires no extra hardware and
reuses frames already being captured.
\end{answersection}

\smartsubsection{Experiment 3-3: Resolution Sweep}
\label{sec:resolution-experiment}

Three target resolutions were benchmarked — $320\times240$, $640\times480$,
and $1280\times720$ — sampling FPS, CPU temperature, and process memory
every five seconds over a five-minute run at each setting, using
\texttt{resolution\_bench.py}.

\smartsubsubsection{An unintentional first result: parameters do not scale for free}

The first pass at $640\times480$ and $1280\times720$ reused the detector
parameters (HOG \texttt{winStride}, \texttt{padding}, \texttt{scale}, and
related settings) that had been tuned for $320\times240$. The result was a
severe, unintended slowdown at the larger resolutions — not a modest
drop, but close to an order of magnitude — since the same
sliding-window/pyramid parameters that were affordable at the smallest
resolution become far more expensive as frame area grows. These
"unoptimized" runs are retained in the results below specifically because
they demonstrate this effect clearly; the corresponding "optimized" rows
show the same resolutions after re-tuning the detection parameters for
each frame size individually (via \texttt{fused\_tracker.py}'s
\texttt{detection\_scale} downsampling, described above).

\begin{table}[H]
\centering
\caption{Resolution sweep results, averaged over a 5-minute run (60 samples, 5s interval) per configuration.}
\label{tab:resolution-sweep}
\resizebox{\textwidth}{!}{%
\begin{tabular}{lcccc}
\toprule
\textbf{Configuration} & \textbf{Avg.\ FPS} & \textbf{Avg.\ Temp.\ (°C)} & \textbf{Max Temp.\ (°C)} & \textbf{Avg.\ Mem.\ (MB)} \\
\midrule
$320\times240$ (baseline) & \stable{1.35} & 52.6 & 55.5 & 78.5 \\
$640\times480$, unoptimized & \alert{0.34} & 54.3 & 57.6 & 92.8 \\
$640\times480$, optimized & \stable{1.61} & 51.2 & 54.2 & 87.9 \\
$1280\times720$, unoptimized & \alert{0.08} & 53.4 & 57.8 & 135.4 \\
$1280\times720$, optimized & \highlight{1.00} & 50.4 & 54.1 & 119.6 \\
\bottomrule
\end{tabular}
}
\end{table}

\begin{figure}[H]
\centering
% \includegraphics[width=0.95\textwidth]{../../2.Image Processing/results/fig-gitignore/320p_1_1_result.png}
\caption{$240\times320$, optimized configuration: annotated frame (left) alongside the live \texttt{resolution\_bench.py} output (right), FPS climbing toward and stabilizing around 1.4.}\label{fig:320p-result}
\end{figure}

\begin{figure}[H]
\centering
% \includegraphics[width=0.95\textwidth]{../../2.Image Processing/results/fig-gitignore/480p_1_1_result.png}
\caption{$640\times480$, unoptimized configuration: annotated frame (left) alongside the live \texttt{resolution\_bench.py} output (right), FPS climbing toward and stabilizing around 0.35.}\label{fig:480p-result}
\end{figure}

\begin{figure}[H]
\centering
% \includegraphics[width=0.95\textwidth]{../../2.Image Processing/results/fig-gitignore/720p_0_7_unoptimized.png}
\caption{$1280\times720$, unoptimized configuration, as we can see the model is overthinking the situation which caused a huge hallucination of the model and costed to count 7 objects while there were no subjects.}\label{fig:720p-unoptimized}
\end{figure}

\begin{figure}[H]
\centering
% \includegraphics[width=0.95\textwidth]{../../2.Image Processing/results/fig-gitignore/720p_1_1_result.png}
\caption{$1280\times720$, optimized configuration: annotated frame (left) alongside the live \texttt{resolution\_bench.py} output (right), FPS climbing toward and stabilizing around 1.0.}\label{fig:720p-result}
\end{figure}

An interesting secondary observation is that the \emph{optimized} runs
show \punder{lower average temperature} than the unoptimized runs at the
same resolution, despite completing more frames per second. This is
consistent with the unoptimized configuration keeping the CPU pinned at
high utilization processing each individual (full-resolution) frame — low
completed-FPS does not mean low CPU load, only that each unit of work
takes far longer. The optimized configuration does genuinely less
per-frame computation (thanks to the pre-detection downsample), so both
FPS improves \emph{and} thermal load drops.

\smartsubsubsection{Optimization Parameters}

Both tuning levers live in \texttt{fused\_tracker.py}: the arguments passed
to the HOG detector's \texttt{detectMultiScale} call (window stride,
padding, pyramid scale, and the grouping threshold), and
\texttt{HOLD\_FRAMES} — how many consecutive frames a track is kept alive
on its cheap visual tracker before the next expensive re-detection is
required. As frame size grows, both were relaxed: a coarser sliding window
and a larger pyramid step reduce the number of detector windows evaluated
per frame, while a higher \texttt{HOLD\_FRAMES} compensates for detection
now running less frequently relative to wall-clock time by letting each
track persist longer between re-detections.

\begin{table}[H]
\centering
\caption{Detection parameters changed between the unoptimized and optimized runs at each resolution.}
\label{tab:optimization-params}
\resizebox{\textwidth}{!}{%
\begin{tabular}{lccc}
\toprule
\textbf{Parameter} & \textbf{$320\times240$ (baseline)} & \textbf{$640\times480$ optimized} & \textbf{$1280\times720$ optimized} \\
\midrule
\texttt{hitThreshold} & 0 & 0 & 0 \\
HOG \texttt{winStride} & $(4,4)$ & $(8,8)$ & $(16,16)$ \\
HOG \texttt{padding} & $(8,8)$ & $(8,8)$ & $(8,8)$ \\
HOG \texttt{scale} & 1.03 & 1.05 & 1.08 \\
\texttt{finalThreshold} (grouping) & 1 & 1 & 2 \\
\texttt{useMeanshiftGrouping} & \texttt{False} & \texttt{False} & \texttt{False} \\
\texttt{HOLD\_FRAMES} & \stable{5} & \stable{7} & \stable{9} \\
\bottomrule
\end{tabular}
}
\end{table}

\smartsubsubsection{Conclusion: Optimal Setting}

Taking FPS, thermal headroom, and memory footprint together, the
$640\times480$ optimized configuration is the strongest overall choice for
continuous deployment: it achieves the \stable{highest average FPS of any
configuration tested} (1.61, even edging out the smaller $320\times240$
baseline) while running at the \stable{lowest average temperature} of the
five configurations. $1280\times720$ optimized remains usable
(1.00~FPS) and would be preferable only if higher spatial detail in the
saved frames is specifically required, at the cost of noticeably higher
memory use and a modest FPS penalty but the fact was that the $320\times240$ was more accurate across the tests, So we decided to go with this one as our best choice.

\smartsubsection{Section Summary}

\begin{table}[H]
\centering
\caption{Image Processing deliverables completed in this section.}
\label{tab:imgproc-summary}
\begin{tabular}{ll}
\toprule
\textbf{Item} & \textbf{Status} \\
\midrule
Camera/video feed (via HTTP-poll workaround) & Complete \\
Lightweight detector (HOG+SVM) & Complete \\
Bounding boxes + live counter overlay & Complete \\
Student ID + live date/time overlay & Complete \\
Live FPS overlay & Complete \\
Shared output for the C layer (\texttt{/dev/shm}) & Complete \\
Exp.\ 3-1 (lighting accuracy) & Complete --- Table~\ref{tab:lighting-accuracy} \\
Exp.\ 3-2 (spoof test) & Complete --- analysis above \\
Exp.\ 3-3 (resolution sweep) & Complete --- Table~\ref{tab:resolution-sweep} \\
\bottomrule
\end{tabular}
\end{table}